% Generated by reproduction/framework_experiments/build_manuscript_assets.py.
\subsection{数据记录、数值溯源与复现}
\label{sx:records}
本补充材料说明第~\ref{sec:analysis}~节所报告实验的完整协议。随附复现目录包含有限元生成器、两个控制脚本、阻尼扫描、新增框架验证脚本及其 JSON 记录。这些标量输入定义了新的数值基准。文献参考性质、制造商范围、研究设计值与计算输出在全文中始终保留各自不同的含义。

记录的运行环境为 Python 3.12.14、NumPy 2.3.5 和 SciPy 1.17.0。
在随附源码包的根目录下，可通过以下命令重新生成完整的数值依赖链：
\begin{verbatim}
cd reproduction
python run_all.py
\end{verbatim}
外层驱动程序先重建有限元基准，再调用实验驱动程序。
后者按以下顺序执行各阶段：
\begin{enumerate}
\item \texttt{minimal\_control\_loop.py};
\item \texttt{route\_b\_shared\_data.py};
\item \texttt{damping\_sensitivity.py};
\item \texttt{summarize\_results.py}.
\end{enumerate}
程序将 OpenBLAS、OpenMP 和 MKL 的线程数均设为一。
路线 B 使用新生成的路线 A 结果进行一致性核查，
同时独立重新生成共享观测并重新拟合系数。

稿件中的数值与科学图由已有 JSON 文件直接生成。从源码包根目录执行
\begin{verbatim}
cd reproduction/framework_experiments
python build_manuscript_assets.py
\end{verbatim}
该写作生成脚本不重新运行实验。其数值溯源 JSON 记录输入哈希、来源字段与显示舍入方式。不同数值库可能使浮点结果末位发生变化；逐位一致不能替代所报告的收敛核查。随附审计还保留实际运行时读取的哈希，并记录随后仅涉及注释的能量说明及历史阻尼元数据修正。该修正没有改变物理参数、模态与敏感性部分。

\subsection{完整物理输入与坐标约定}
\label{sx:physical-inputs}
物理模型是一根均匀、实心、平面的 Euler--Bernoulli 梁，
根部固支，自由端带有理想刚性磁芯。
其自由长度为 $L=0.08$ m，半径为 $r=0.0006$ m，
近似杨氏模量为 $E=2.1\times10^6$ Pa，近似密度为 $\rho=1028$ kg/m$^3$。
后两个数值是一组配对材料参考值，对应约 $10$ wt\% 的固化剂质量分数，
以及在 $80\,^{\circ}$C 下固化 $2$ h 的条件 \citep{cho2021pdms}。
固化剂占完整混合物的 $10$ wt\%，约对应 $9{:}1$ 的基料与固化剂质量比。
这些数据并不意味着本文的目标圆梁就是该研究测量的试样。

裸磁芯设计为立方体，边长 $a=0.00025$ m，
标称密度 $\rho_m=7500$ kg/m$^3$，剩磁 $B_r=1.425$ T。
用于比较的制造商范围为 $1.40$--$1.45$ T，
密度取其声明的下限 \citep{smN50}。
有量纲磁偶极矩估计为 $\mu\simeq B_ra^3/\mu_0$，其中
$\mu_0\simeq4\pi\times10^{-7}$ H/m \citep{kjMoment}。
这是针对均匀磁化硬磁体的工程近似。
假定磁芯中心与梁端点重合，磁化方向沿直参考梁的切向。
恒定纵向磁场为 $B_0=0.03$ T。
模型不包含横向重力、接触、外部流动、梁的分布转动惯量、
剪切变形、镀层、胶黏剂或夹具质量。
每一项省略都是建模假设。

\begin{table}[t]
\centering
\caption{由标称输入推导的物性量。数值取自
\texttt{new\_baseline.json}；显示精度描述计算精度，而非实验精度。}
\label{sx:properties}
\setlength{\tabcolsep}{3pt}
\begin{tabular}{lcl}
\toprule
物理量 & 计算值 & 单位\\
\midrule
截面积 $A$ & $1.13097336\times10^{-6}$ & m$^2$ \\
截面二次矩 $I$ & $1.01787602\times10^{-13}$ & m$^4$ \\
抗弯刚度 $EI$ & $2.13753964\times10^{-7}$ & N\,m$^2$ \\
梁质量 $\rho AL$ & $9.30112487\times10^{-5}$ & kg \\
磁芯体积 $V_m$ & $1.56250000\times10^{-11}$ & m$^3$ \\
磁芯质量 $m_t$ & $1.17187500\times10^{-7}$ & kg \\
磁芯转动惯量 $J_t$ & $1.22070312\times10^{-15}$ & kg\,m$^2$ \\
磁偶极矩 $\mu$ & $1.77184214\times10^{-5}$ & A\,m$^2$ \\
能量尺度 $E_0$ & $5.31552642\times10^{-7}$ & J \\
\bottomrule
\end{tabular}
\end{table}

有量纲模态基满足 $\Phi^{\mathsf T}M_{\mathrm{total}}\Phi=I$。
其平衡位移尺度为 $S_j=\sqrt{E_0/(\omega_0\omega_j)}$，
其中 $E_0=\mu B_0$ 且 $\omega_0=\omega_1$。
因此，物理模态位移与动量分别为 $a_j=S_jX_j$ 和
$\pi_j=\sqrt{E_0\omega_j/\omega_0}\,Y_j$。
在无量纲时间 $\tau=\omega_0t$ 下，
$X'=\mathsf NY$，二阶 $X$ 方程中的等效质量为 $\mathsf N^{-1}$。
平衡缩放后，该质量矩阵不是单位矩阵。
端点转角与位移分别为 $d^{\mathsf T}X$ 和 $\ell^{\mathsf T}X$，
其中系数见式~\eqref{num:scales}。
全部物理误差指标均使用同一映射。

\subsection{有限元与独立边界核查}
\label{sx:fe}
每个三次 Hermite 单元采用通常的一致平动质量与弯曲刚度。
消去根部位移与斜率后，保留 $2n_e$ 个位形坐标。
磁芯在末端位移质量项上增加 $m_t$，
在末端转动惯量项上增加 $J_t=m_ta^2/6$。
偏置磁场在末端转动刚度项上增加 $\mu B_0$。
特征值求解采用质量归一化，并以缩放后的逆算子构造对称特征值问题，
直接分辨最低阶模态，避免从高频谱端作相减运算来恢复它们。
前四阶频率用于网格核查，前两阶用于定义控制模型。

独立连续体核查求解 $EIw''''=\rho A\omega^2w$，
边界条件为 $w(0)=w'(0)=0$ 及
\begin{equation}
 \begin{split}
 EIw'''(L)+m_t\omega^2w(L)&=0,\\
 EIw''(L)+(\mu B_0-J_t\omega^2)w'(L)&=0.
 \end{split}
 \label{sx:boundary-check}
\end{equation}
满足固支条件的基函数为 $\cosh(kx)-\cos(kx)$ 和 $\sinh(kx)-\sin(kx)$。
对其端点行列式独立确定根的括区间，
再采用绝对求根容差为 $2\times10^{-13}$ 的 Brent 迭代求解。
这针对同一理想边界模型给出独立浮点核查，不是特征值的区间证明。

\begin{table*}[t]
\centering
\caption{有限元频率（rad/s），以及前四阶模态相对于独立连续体根的最大相对误差。
采用 $32$ 个单元的结果定义本文基准。}
\label{sx:fe-table}
\setlength{\tabcolsep}{5pt}
\begin{tabular}{rccccc}
\toprule
单元数 & $\omega_1$ & $\omega_2$ & $\omega_3$ & $\omega_4$ & 最大相对误差\\
\midrule
16 & $7.84001022$ & $47.39980944$ & $131.21158140$ & $256.37539514$ & $1.5071\times10^{-4}$ \\
32 & $7.84000916$ & $47.39957662$ & $131.20668107$ & $256.33922193$ & $9.5903\times10^{-6}$ \\
64 & $7.84000910$ & $47.39956199$ & $131.20637133$ & $256.33691793$ & $6.0208\times10^{-7}$ \\
\bottomrule
\end{tabular}
\end{table*}

表~\ref{sx:sensitivity-table} 每次仅改变一个声明的输入。
磁芯质量系数同时乘于 $m_t$ 和 $J_t$，并保持几何与磁偶极矩不变。
这是独立的惯性敏感性分析，不是制造商公差或完全耦合的制造模型。
特别是，系数 $0.9$ 不被解释为所采用密度规格内的同尺寸磁芯。
$20\%$ 的模量变化同样是设计扰动。
JSON 中还保留了剩磁与惯性组合的四个角点。
在这些有限元参数敏感性分析中，标称有量纲 $\alpha$ 始终固定，
因此所得到的阻尼比随频率变化。
这一约定与另行进行的阻尼扫描有意区分：后者保持物理基准不变而改变 $\alpha$。

\begin{table*}[t]
\centering
\caption{选定的有限元设计敏感性。末列采用固定的标称
$\alpha=0.3920004581$ s$^{-1}$。这些数值均非实测置信区间。}
\label{sx:sensitivity-table}
\begin{tabular}{lccc}
\toprule
扰动 & $\omega_1$（rad/s） & 变化（\%） & $\alpha/(2\omega_1)$\\
\midrule
$B_r=1.40$ T & $7.83350702$ & $-0.08294$ & $0.02502075$ \\
$B_r=1.45$ T & $7.84648941$ & $+0.08266$ & $0.02497935$ \\
$m_t,J_t$ 乘以 $0.9$ & $7.84192087$ & $+0.02438$ & $0.02499391$ \\
$m_t,J_t$ 乘以 $1.1$ & $7.83809879$ & $-0.02437$ & $0.02500609$ \\
$E$ 乘以 $0.8$ & $7.09311163$ & $-9.52674$ & $0.02763247$ \\
$E$ 乘以 $1.2$ & $8.51964699$ & $+8.66884$ & $0.02300567$ \\
\bottomrule
\end{tabular}
\end{table*}

\subsection{训练数据、验证数据与积分拟合}
\label{sx:training}
标称合成受控对象取 $\beta_*=0.2$ 和 $\gamma_*=0.05$，并保留完整余弦。
训练状态和验证状态由经典 RK4 以步长 $0.0005$ 在 $0\leq\tau\leq4$ 上生成。
每条记录的物理时长为 $0.510203$ s，并非四秒。
每隔四个生成步长采样一次，可得到每条轨迹的 $2001$ 个状态，观测间隔为 $0.002$。
不施加随机噪声、平滑或导数估计。
六个固定初始状态列于表~\ref{sx:initial-states}。
它们针对这一更柔软的梁重新设计，采用更小振幅，
每个状态均为前一个最小实验中对应初始状态的一半。
它们不是从实测分布中取得的样本。

\begin{table}[t]
\centering
\caption{无量纲初始状态。验证轨迹不参与参数拟合。}
\label{sx:initial-states}
\setlength{\tabcolsep}{3.3pt}
\begin{tabular}{lrrrr}
\toprule
记录 & $X_1$ & $X_2$ & $Y_1$ & $Y_2$\\
\midrule
训练 1 & $0.09$ & $0.07$ & $0.025$ & $-0.02$ \\
训练 2 & $-0.075$ & $0.1$ & $0.01$ & $0.025$ \\
训练 3 & $0.11$ & $-0.06$ & $-0.015$ & $0.03$ \\
训练 4 & $0.05$ & $0.12$ & $0.02$ & $-0.01$ \\
验证 1 & $-0.095$ & $0.055$ & $0.03$ & $0.005$ \\
验证 2 & $0.065$ & $-0.105$ & $-0.02$ & $0.015$ \\
\bottomrule
\end{tabular}
\end{table}

每个互不重叠时窗包含 $20$ 个区间和 $21$ 个状态。
相邻时窗共享边界观测。
复合 Simpson 权重直接作用于式~\eqref{num:regression} 中的两个积分。
每条轨迹提供 $100$ 个向量时窗：四条训练轨迹给出 $800$ 个标量行，
两条验证轨迹给出 $400$ 个标量行。
将两列设计向量分别除以各自的 Euclidean 范数，
采用无约束最小二乘例程求解，再恢复原尺度。
不使用正则化或非负约束。
源脚本检验了矩阵满列秩。

归一化联合 Gram 矩阵 $G=A^{\mathsf T}A/800$ 为
\begin{equation}
 G=\begin{pmatrix}5.16502121\times10^{-14} & 2.38379831\times10^{-11}\\2.38379831\times10^{-11} & 5.59508836\times10^{-6}\end{pmatrix}.
 \label{sx:gram}
\end{equation}
这里按标量行数归一化；若改用向量时窗数，系数界中的数值常数会发生变化。
令 $r_I=\|A\widehat\theta-b\|_2/\sqrt{800}$，
若能独立确立缺陷界 $\|A\theta_*-b\|_2/\sqrt{800}\leq\varepsilon_I$，
即可得到第~\ref{sec:kam}~节描述的条件性估计。
利用已知合成真值进行诊断，得到
$\varepsilon_{I,\mathrm{diag}}=1.26876\times10^{-12}$，相应界表达式的数值为
$1.00470\times10^{-5}$，而实际系数误差为 $3.36849\times10^{-6}$。
该诊断使用了仅在合成实验中可知的真值，以及普通浮点求积。
它不是 $\varepsilon_I$ 的独立包围，不能据此完成严格的解析范数或 KAM 论证。
留出验证残差考察同一特征族内的不同轨迹，
并未检验更大的特征类或含噪观测模型。

\subsection{路线 A 协议与中点法验证}
\label{sx:route-a}
初始状态为 $(0.09,0.08,0.015,-0.01)$，目标为 $h_0$ 的保守流。
评估以下四种指令：
$U=0$、$U=\widehat\gamma Y$、
$U=\widehat\gamma Y+\widehat\beta\phi(X)$，以及
$U=\gamma_*Y+\beta_*\phi(X)$。
最后一种为真值参数对照。
各指令均通过相同的理想物理端口例程重构，
重构值随后代入受控对象的微分方程。
因此，端口逆映射在积分过程中实际参与计算，
而非仅在模态仿真完成后进行评估。

主要中点步长为 $0.001$，终止时间为 $\tau=20$，
保存全部端点，共得到 $20001$ 个状态。
物理时间步长为 $1.27551\times10^{-4}$ s；
这里不声称这一数值是硬件能够实现的采样周期。
每次 Newton 求解均由显式 Euler 预测值初始化。
停止准则为
\begin{equation}
 \left\|x_{k+1}-x_k-
 \Delta f\!\left(\frac{x_k+x_{k+1}}2\right)\right\|_\infty
 \leq2\times10^{-14}.
 \label{sx:newton-stop}
\end{equation}
迭代使用解析中点 Jacobian，连续 $12$ 次修正尝试仍未成功时抛出错误。
此准则针对完整增量方程，而非仅针对修正量范数。
记录的路线 A 主要计算最多使用一次修正；较粗步长的收敛阶核查最多使用两次。

独立 $h_0$ 参考由 RK4 计算，步长分别为 $0.0005$ 和 $0.00025$。
主要比较每隔四步对细步长参考采样一次。
真值参数对照的中点收敛测试采用步长 $0.04$、$0.02$、$0.01$ 和 $0.005$，
始终在共同时间点与上述独立细步长参考进行比较。

\begin{table}[t]
\centering
\caption{路线 A 真值参数对照的中点法收敛，终止时间为 $\tau=20$。
最大值取遍所有保存的共同时间点。}
\label{sx:order-table}
\setlength{\tabcolsep}{4pt}
\begin{tabular}{rcc}
\toprule
$\Delta$ & 最大状态误差 & 观测收敛阶\\
\midrule
$0.04$ & $4.64310\times10^{-2}$ & --- \\
$0.02$ & $1.18409\times10^{-2}$ & $1.971306$ \\
$0.01$ & $2.96760\times10^{-3}$ & $1.996414$ \\
$0.005$ & $7.42245\times10^{-4}$ & $1.999331$ \\
\bottomrule
\end{tabular}
\end{table}

观测收敛阶为 $\log_2(E_\Delta/E_{\Delta/2})$。
辛性诊断采用 $D\Psi^{\mathsf T}JD\Psi-J$ 的谱范数，
其中 $D\Psi$ 是在求得的中点处计算的理想隐式 Cayley 切映射。
仅对真值参数对照恢复使用这一诊断，以检验哈密顿目标映射。
随附的矩阵消元核查另外采用 $n=1,2,4$ 个位形坐标，
验证二阶更新的代数关系；它不是高阶模态梁精度实验。

\subsection{重新生成观测的阻尼扫描}
\label{sx:damping}
五个第一阶模态阻尼比分别为 $0$、$0.01$、$0.025$、$0.05$ 和 $0.1$。
对于每个值，有量纲阻尼系数为 $\alpha=2\zeta_1^{\mathrm{d}}\omega_0$，
真实无量纲阻尼系数为 $\gamma=2\zeta_1^{\mathrm{d}}$。
有限元基底、人工设定的 $\beta_*=0.2$ 及全部初始状态保持不变。
重新生成所有训练与验证观测，
并采用不变的无约束双列拟合过程重新拟合两个系数。
因此，没有任何数值是通过复制基准拟合结果或仅改变控制器参数得到的。

\begin{table*}[t]
\centering
\caption{每个阻尼情景均重新拟合。零阻尼情景中的极小负估计值来自数值计算，未施加约束将其消除。}
\label{sx:damping-fit}
\setlength{\tabcolsep}{5pt}
\begin{tabular}{rcccc}
\toprule
$\zeta_1^{\mathrm{d}}$ & $\alpha$（s$^{-1}$） & $\widehat\beta$ & $\widehat\gamma$ & $\lambda_{\min}(G)$\\
\midrule
$0$ & $0.000000000$ & $0.200003040949$ & $-1.4037896716\times10^{-12}$ & $6.9451\times10^{-14}$ \\
$0.01$ & $0.156800183$ & $0.200003177492$ & $1.9999999994\times10^{-2}$ & $6.1346\times10^{-14}$ \\
$0.025$ & $0.392000458$ & $0.200003368491$ & $4.9999999987\times10^{-2}$ & $5.1549\times10^{-14}$ \\
$0.05$ & $0.784000916$ & $0.200003636621$ & $9.9999999975\times10^{-2}$ & $3.9870\times10^{-14}$ \\
$0.1$ & $1.568001832$ & $0.200003938495$ & $1.9999999995\times10^{-1}$ & $2.6883\times10^{-14}$ \\
\bottomrule
\end{tabular}
\end{table*}

恢复路线采用中点步长 $0.002$，积分至 $\tau=10$，即 $1.275509$ s。独立 $h_0$ 参考采用 RK4 步长 $0.00025$。包含无控制、学习及真值参数对照组；主要实验中的仅阻尼补偿组不在此重复。这里所有误差均使用完整 $[0,10]$ 时窗。补充跟踪扩展使用同一组重新拟合的参数，但采用自身的目标与后半时窗；这些不同结果保留在表~\ref{sx:damping-errors}。主要实验的收敛核查不能代替对每个扫描点单独开展收敛研究。

\begin{table*}[t]
\centering
\caption{方向一的阻尼情景，每个情景均重新生成观测并重新拟合。所有数值均为相对独立 $h_0$ 参考的状态 RMS，评价时窗为 $[0,10]$。}
\label{sx:framework-damping-a}
\begin{tabular}{rccc}
\toprule
$\zeta_1^{\mathrm d}$ & 无控制 & 学习恢复 & 真值参数恢复\\
\midrule
$0$ & $3.77362\times10^{-5}$ & $3.42885\times10^{-5}$ & $3.42885\times10^{-5}$ \\
$0.01$ & $6.73657\times10^{-3}$ & $3.42885\times10^{-5}$ & $3.42885\times10^{-5}$ \\
$0.025$ & $1.59501\times10^{-2}$ & $3.42885\times10^{-5}$ & $3.42885\times10^{-5}$ \\
$0.05$ & $2.92396\times10^{-2}$ & $3.42885\times10^{-5}$ & $3.42885\times10^{-5}$ \\
$0.1$ & $4.97374\times10^{-2}$ & $3.42885\times10^{-5}$ & $3.42885\times10^{-5}$ \\
\bottomrule
\end{tabular}
\end{table*}

\subsection{指标定义、物理端口与采样极值}
\label{sx:metrics}
对于保存的误差 $e_k=x_k-x_{\mathrm{ref}}(\tau_k)$、$k=0,\ldots,n$，状态指标为
\begin{equation}
 E_{\mathrm{RMS}}=
 \left(\frac1{n+1}\sum_{k=0}^{n}\|e_k\|_2^2\right)^{1/2}.
 \label{sx:rms}
\end{equation}
这是四个平衡无量纲坐标中的欧氏误差，不再额外除以四。后半时窗 RMS 对指定保存状态子集采用同一定义。最大状态误差为 $\max_k\|e_k\|_2$，终点误差为 $\|e_n\|_2$。物理末端位移误差先由 $\ell^{\mathsf T}e_{X,k}$ 计算，末端角度误差先由 $d^{\mathsf T}e_{X,k}$ 计算，再取标量 RMS 或最大值。附加跟踪实验在补充材料第~\ref{sx:route-b}~节中保留独立参考与同步参考各自的指标。

路线 A 的能量漂移为 $\max_k|h_0(x_k)-h_0(x_0)|$，
其相对形式再除以 $|h_0(x_0)|$。
这是标称能量，并非精确的非线性作用量，也不是所有扰动受控对象各自的能量。
指令代价采用 $\int\|U\|_2^2\,\mathrm d\tau$ 的梯形求积近似；
它是无量纲广义控制代价，而非执行器电功率。
回归残差是另一种指标，即 $A\widehat\theta-b$ 各标量行的均方根误差。

局部端口重构假定力与力矩可独立指令，并具有精确状态信息。
将重构磁场与梯度代回模态方程，记录广义力差值。
局部总磁场为 $\sqrt{B_0^2+B_y^2}$，
局部梯度大小为其对称、无迹二维矩阵的谱范数。
没有实现执行器电流映射或饱和限制。

\begin{table*}[t]
\centering
\caption{主要学习恢复实验中的采样磁需求。计入端点和积分中点。这些是理想局部需求，并非实测执行器能力。}
\label{sx:port-table}
\begin{tabular}{lcccc}
\toprule
实验 & $\max|B_y|$ (mT) & $\max|B|$ (mT) & $\max\|\nabla B\|_2$ (T/m) & $\max|F_y|$ ($\mu$N)\\
\midrule
A：学习 & $0.544490$ & $30.004941$ & $0.01603429$ & $0.284102$ \\
\bottomrule
\end{tabular}
\end{table*}

轨迹的角度与位移最大值在保存时间点处采样。
全梁斜率与位移在 $257$ 点空间网格上通过 Hermite 插值计算，
每个单元划分为八个子区间。
随附节点模态形状，因此可另行检查空间多项式极值；
但所报告数值使用上述网格，不包围连续时间极值。
指令诊断包含中点积分器使用的所有端点与算术中点。
作为设计比较讨论的 $20$ mT 横向场值并非强制截幅规则。

\subsection{P1 框架矩阵验证}
\label{sx:framework-matrix}
\texttt{reproduction/framework\_experiments} 单独保存附加实验。从源码包根目录执行的完整入口为
\begin{verbatim}
cd reproduction/framework_experiments
python run_all.py
\end{verbatim}
可选参数 \texttt{--experiment} 用于选择 \texttt{matrix}、\texttt{closed}、\texttt{defect} 或 \texttt{weighted}。所有参数路径均相对脚本解析。已有有限元、训练和恢复记录仅被读取，不会被覆盖。新增 JSON 文件包含基准与源代码哈希以及数值库版本。随附写作生成脚本生成当前三份数值章节输入，其中包括补充跟踪扩展；旧双路线写作模板仅作历史存档。

对于式~\eqref{num:matrix-data} 中的无量纲矩阵族，按 $(q,p)$ 排列的初态为
\begin{align*}
 n=1:&\quad (0.1,0),\\
 n=2:&\quad (0.1,0.06,0,0.02),\\
 n=4:&\quad (0.1,0.06,0.04,0.025,0,0.02,-0.01,0.005).
\end{align*}
每个维数均采用 $T=4$ 和四个步长 $0.08,0.04,0.02,0.01$。SciPy 矩阵指数独立计算 $\exp(tA)x_0$，覆盖最细网格的全部 $401$ 个点。每个较粗步长的方法在这些时刻的对应子集上比较。普通映射与其消元形式独立推进；对每个维数与步长保存最大轨迹差异，同时记录 $I+E_\Delta$ 的最小奇异值。

加权映射为 $S_{1/2}\operatorname{Mid}_\Delta(A_0)S_{1/2}$，其中 $S_{1/2}=\operatorname{diag}(I_n,e^{-0.2\Delta/2}I_n)$。保守核查采用 $K+K_s$，其中 $K_s=0.05\operatorname{diag}(1,\ldots,n)$，对应声明输入 $U=0.2p-K_sq$。对于线性保守向量场 $Ax$，Central4 采用增量矩阵 $G_\Delta=\Delta A-\Delta^3A^3/12$ 与更新 $(I-G_\Delta/2)^{-1}(I+G_\Delta/2)$，并通过线性方程求解实现。附加的 $C_g$ 族只核查消元；非比例阻尼出现在 $n=2,4$，标量成员只用作对照。

\begin{table*}[t]
\centering
\caption{P1 矩阵在 $\Delta=0.01$ 时的全网格最大误差。各列采用各自声明的含阻尼或保守目标。全部更粗步长的误差与观测阶均保留在 JSON 中。}
\label{sx:framework-matrix-errors}
\begin{tabular}{rcccc}
\toprule
$n$ & 普通中点 & 加权中点 & 保守中点 & 保守 Central4\\
\midrule
1 & $1.98924\times10^{-6}$ & $1.97226\times10^{-6}$ & $3.41418\times10^{-6}$ & $3.35319\times10^{-11}$ \\
2 & $3.71730\times10^{-6}$ & $3.85870\times10^{-6}$ & $5.58365\times10^{-6}$ & $8.05545\times10^{-11}$ \\
4 & $6.08098\times10^{-6}$ & $5.58130\times10^{-6}$ & $9.62132\times10^{-6}$ & $1.95812\times10^{-10}$ \\
\bottomrule
\end{tabular}
\end{table*}

\subsection{P3 固定误差闭环向量场与保守 Central4}
\label{sx:framework-closed}
新增磁模型实验读取同一份 \texttt{new\_baseline.json}，使用初态 $(0.09,0.08,0.015,-0.01)$，积分区间为 $\tau\in[0,20]$。固定误差实验设定 $\delta\beta=0.02$、$\delta\gamma=0.01$，受控对象参数为 $\beta_*=0.2$、$\gamma_*=0.05$，因此声明的控制器系数为 $0.18$ 和 $0.04$。这些误差是有意人为设置的，不作为新的辨识结果。步长改变时不重新拟合任何系数。每次中点求解内部都调用已有物理端口重构。独立参考直接实现约化向量场式~\eqref{num:fixed-error-field}，因此没有复用中点算法。

独立 DOP853 参考的较粗设置为 $\mathrm{rtol}=10^{-11}$、$\mathrm{atol}=10^{-13}$、最大步长 $0.05$；精细设置为 $\mathrm{rtol}=2\times10^{-13}$、$\mathrm{atol}=2\times10^{-15}$、最大步长 $0.025$。在相同保存时刻评价稠密输出，并核查两次求解的成功标志。这些是浮点加密核查，并非解的严格包络。中点步长为 $0.04,0.02,0.01,0.005,0.0025$。实际轨迹、标称参考轨迹与闭环参考轨迹的完整数据保存于 \texttt{closed\_loop\_convergence\_traces.npz}。

非线性 Central4 实验固定采用保守 $h_0$（$\kappa=1$，残余阻尼为零），并使用上述前四个步长。它实现式~\eqref{ss:central4-magnetic}--\eqref{ss:central4-magnetic-jacobian} 中的完整增量和解析 Jacobian。复步长核查采用 $10^{-25}i$ 扰动，在主要初态及 $(-0.07,0.04,0.03,-0.02)$ 处进行，步长为 $\Delta=0.04$。最大导数差异为 $2.7756\times10^{-17}$。该实验不声称不完全补偿的物理闭环具有哈密顿结构，也不声称追加反馈跳变仍保持四阶。

全部隐式框架求解均采用增量预测值、残差无穷范数容差 $2\times10^{-14}$，以及至多 $12$ 次 Newton 修正；失败即停止运行。对于中点法，预测值为显式 Euler；对于 Central4，预测值由前一状态处的完整增量给出。每次运行都记录最大实际残差、最大修正次数、失败步数，以及采样 Newton 矩阵的最小奇异值。标准与指数加权缺陷评价的是计算中点处的理想隐式切映射，而非有限次浮点迭代的导数。式~\eqref{ctrl:global-error} 的一般范数结论在统一使用相容范数时成立；对磁模型的四维状态，无穷范数残差给出的欧氏残差上界至多为其两倍。

\subsection{人为构造的代数缺陷预算}
\label{sx:framework-defect}
这一独立单自由度实验采用
\begin{equation}
 A=\begin{pmatrix}0&1\\-1&0\end{pmatrix},\quad
 x_0=(0.1,0.02)^{\mathsf T},\quad T=4,
 \label{sx:defect-model}
\end{equation}
以及步长 $0.08,0.04,0.02,0.01$。人为求解
\begin{equation}
 \begin{aligned}
 x_{k+1}-x_k-\Delta A\frac{x_k+x_{k+1}}2
   &=\eta_\Delta(1,0)^{\mathsf T},\\
 \eta_\Delta&=0.01\Delta^q,
 \end{aligned}
 \label{sx:defect-injection}
\end{equation}
其中 $q=1,2,3$。更新等于普通中点矩阵作用的结果加上 $(I-\Delta A/2)^{-1}\eta_\Delta(1,0)^{\mathsf T}$。因此，残差被有意引入代数方程，而非归因于 Newton 求解器。JSON 保存每个实际残差向量及其无穷范数、最大欧氏范数，以及相对预设残差的最大差异。整个实验中，最后一项为 $3.5779\times10^{-17}$。

保留两种不同误差：相对 $\exp(tA)x_0$ 的最大总误差，以及相对相同步长下独立推进的未扰动中点轨迹的最大新增误差。新增误差单独反映 $\eta_\Delta/\Delta$ 项；总误差还包含普通的 $O(\Delta^2)$ 积分项。

\begin{table*}[t]
\centering
\caption{$\Delta=0.01$ 时的人为缺陷诊断。最终观测阶由步长 $0.02$ 与 $0.01$ 比较得到。“新增”指相对未扰动中点法的误差，“总量”指相对连续振子的误差。}
\label{sx:defect-table}
\begin{tabular}{rccccc}
\toprule
$q$ & 实际残差最大 $\infty$ 范数 & 总误差 & 总误差阶 & 新增误差阶 & 新增误差$/ (\eta_\Delta/\Delta)$\\
\midrule
1 & $1.000\times10^{-4}$ & $2.00026\times10^{-2}$ & $0.000566$ & $-0.000000$ & $1.999999$ \\
2 & $1.000\times10^{-6}$ & $2.02630\times10^{-4}$ & $1.018776$ & $1.000000$ & $1.999999$ \\
3 & $1.000\times10^{-8}$ & $5.18612\times10^{-6}$ & $1.999969$ & $2.000000$ & $1.999999$ \\
\bottomrule
\end{tabular}
\end{table*}

\subsection{非线性指数加权对照}
\label{sx:framework-weighted}
加权实验使用两个明确已知的向量场：未补偿合成受控对象 $(\delta\beta,\delta\gamma)=(0.2,0.05)$，以及固定误差诊断 $(0.02,0.01)$。它采用与主要磁模型实验相同的初态、$T=4$ 和步长 $0.04,0.02,0.01,0.005$。对于每个向量场，普通中点法与加权映射均同上述设置的独立 DOP853 参考比较。加权实现为 $S_{1/2}\operatorname{Mid}_{h_\kappa}S_{1/2}$，其中 $S_{1/2}=\operatorname{diag}(I_2,e^{-\delta\gamma\Delta/2}I_2)$，$\kappa=1+\delta\beta$。所得数值解还代回式~\eqref{ctrl:weighted-midpoint} 核查。加权理想切映射为 $S_{1/2}D\operatorname{Mid}_{h_\kappa}S_{1/2}$，故目标因子是 $e^{-\delta\gamma\Delta}$。两个模型及全部误差分别记录；实际部署的学习控制器没有读取未知真实误差。

\begin{table*}[t]
\centering
\caption{非线性加权映射诊断。缺陷为 $\Delta=0.04$ 时指数加权谱范数缺陷的最大值；观测阶使用最后一次步长减半。每组比较中的两种方法逼近同一向量场。}
\label{num:nonlinear-weighted}
\begin{tabular}{lcccc}
\toprule
向量场 $(\delta\beta,\delta\gamma)$ & 普通缺陷 & 加权缺陷 & 普通观测阶 & 加权观测阶\\
\midrule
$(0.2,0.05)$ & $2.8763\times10^{-5}$ & $3.3771\times10^{-16}$ & $1.999452$ & $1.999400$ \\
$(0.02,0.01)$ & $5.7618\times10^{-6}$ & $3.3603\times10^{-16}$ & $1.999414$ & $1.999403$ \\
\bottomrule
\end{tabular}
\end{table*}

\subsection{条件性净空与历史结果的迁移限制}
\label{sx:migration}
$4$ mm 净空是扣除杆半径后的中心线净余量。
标称能量比较采用 $E_c=3EIc^2/(2L^3)=10.019717$ nJ，
并选取系数 $0.95$。
它通过非负的标称动能、弹性能与静磁场能定义了条件性充分界。
主要实验的初始能量为 $12.656721$ nJ，超过所提出的 $9.518731$ nJ 能量上限。
当前控制没有实现输入功约束器。
因此，采样位移小于净空不能被报告为已验证该能量上限前提。

原稿中的角度、角速度、速度和执行器误差预算，原本就已声明为确定性工程假设。
它们仍保持这种假设地位，不能重新归类为实测传感器规格。
在将其用于保证之前，必须针对新的物理映射，
重新计算运行裕量及任何延迟界。

\begin{table*}[t]
\centering
\caption{改变材料输入、端点惯性与模态映射后，主要数值结论的处理方式。
完整的章节级登记表随源码包提供。}
\label{sx:migration-table}
\begin{tabular}{p{.23\textwidth}p{.29\textwidth}p{.41\textwidth}}
\toprule
原有组成部分 & 当前处理 & 建立新数值结论所需工作\\
\midrule
梁能量、复数表示与控制恒等式 & 使用新的质量与坐标映射保留推导 & 针对所选残差与端口，核查全部特化假设及导数。\\
有限元频率与平衡映射 & 已在 $16/32/64$ 个单元下重新计算 & 高阶模态受迫轨迹仍需单独验证。\\
匹配条件下的实数／复数 CHNN 拟合及训练随机种子对比 & 作为历史结果；本文用明确声明的双标量测试替代 & 生成新数据，执行完整的匹配训练与验证协议。\\
GC17 中点法／Central4 表 & 不采用历史数值 & 重新拟合目标能量，并在该函数上重复两种积分测试。\\
COMSOL 全有限元、两模态与模型失配比较 & 不采用历史数值 & 重建每个独立模型、输入、求解器协议及降阶比较。\\
有限振幅 KAM、族／测度与尾项证书 & 保留条件性理论；不采用旧数值证书 & 重新计算拟合函数、解析域、Gram／缺陷界、扭转、环面及经验证的余项不等式。\\
全梁环面提升与持续控制界 & 保留条件性构造；没有新的全阶仿真 & 重新计算高阶模态、尾项强迫、输入界及启动暂态。\\
输入功约束器、运行容差与 S 形路径结果 & 保留条件性论证；不采用旧输出 & 针对新模型重新初始化并实现相应控制器及裕量。\\
长期与共振边界诊断 & 不采用历史输出 & 重新执行所述时域、采样、延拓及收敛研究。\\
\bottomrule
\end{tabular}
\end{table*}

特别地，旧有 $22.66$ nJ 能量预算不是新基准输入。
旧环面频率、Fourier 系数、$35$ 参数拟合、COMSOL 轨迹、
$200$ s 预测统计及全阶跟踪误差均描述不同的数值对象。
它们的来源可以保留在历史登记表中，但不作为本文所报告计算的输出。
因此，重新计算的两模态测试与保留的条件性定理，具有明确不同的证据作用。
